\begin{table}
\caption{Konfound (Frank et al., 2013): pkonfound para HLM e OLS por modelo. Δ pkonfound mostra a subestimação sistemática da robustez pelo OLS quando aplicado a dados com estrutura multinível (ICC $>$ 5\%). PNAD Contínua 2016--2025.}
\label{tab:konfound_hlm_vs_ols}
\begin{tabular}{lrrrrrrrr}
\toprule
Modelo & β_negro & SE (HLM) & t (HLM) & pkonfound HLM (%) & SE (OLS) & t (OLS) & pkonfound OLS (%) & Δ pkonfound (pp) \\
\midrule
M1 & -0.214900 & 0.001300 & -165.300000 & 98.800000 & 0.009400 & -22.900000 & 91.400000 & 7.400000 \\
M2 & -0.102000 & 0.001300 & -78.500000 & 97.500000 & 0.003900 & -26.200000 & 92.500000 & 5.000000 \\
M3 & -0.102000 & 0.001300 & -78.500000 & 97.500000 & 0.003900 & -26.200000 & 92.500000 & 5.000000 \\
M4 & -0.063700 & 0.001200 & -53.100000 & 96.300000 & 0.004300 & -14.800000 & 86.800000 & 9.500000 \\
\bottomrule
\end{tabular}
\end{table}
